\documentclass[11pt]{article}
\usepackage[margin=1in]{geometry}
\usepackage{hyperref}
\usepackage{xcolor}
\usepackage{amsmath}
\usepackage{booktabs}
\usepackage{enumitem}

\newcommand{\editortext}[1]{\textit{\color{gray}#1}}
\newcommand{\newtext}[1]{\begin{quote}\small #1\end{quote}}

\title{Response to the Editor\\[0.5em]
\large Revised Manuscript: ``Evaluating Tool-Using Language Agents:\\
Judge Reliability, Propagation Cascades, and Runtime Mitigation\\
in AgentProp-Bench''}
\author{Anonymous (for review)}
\date{\today}

\begin{document}
\maketitle

\section{Overview of Changes}

We thank the editor for the constructive feedback. The manuscript has been
substantially revised:

\begin{itemize}[nosep]
  \item Renamed from AgentHallu-Bench to \textbf{AgentProp-Bench} (see Section~9 below).
  \item \textbf{Repositioned} around evaluation reliability as the primary
        contribution, with propagation and mitigation as evidence.
  \item Added 18 new 2024--2026 references from IR/IS and agent-hallucination venues.
  \item Added standalone sections for Research Objectives (Section~3),
        Dataset (Section~4), and Discussion (Section~7).
  \item Expanded from $\sim$5,200 to $\sim$9,500 words.
  \item Corrected a methodological issue in stage-indicator definitions
        (disclosed transparently in Section~9).
  \item Collected 100 stratified human labels to validate the ensemble judge.
  \item Re-ran the interceptor under a concurrent no-interceptor control.
\end{itemize}

\section{Response to Editor Point 1: Abstract}

\editortext{``Your abstract does not highlight the specifics of your
research or findings but contains too much background information.
Some details of your research would be nice, for example, numbers
addressing the sample, data, percentage improvement, etc.''}

\textbf{Addressed.} The abstract has been rewritten to 195 words, leading
with specific numbers rather than background. It now includes:
sample sizes (2,000 tasks, 2,300 traces, 9 LLMs, 100 human labels),
key results ($\kappa=0.049$ for heuristic, $\kappa=0.432$ for ensemble,
$r_{2,3}\approx 0.62$ human-calibrated propagation, $-23$\,pp interceptor
effect), and methodology (concurrent $n{=}600$ control). Background material
was reduced to one sentence.

\newtext{See the revised Abstract (page~1). Word count: ${\sim}$170 (down from 247).}

\section{Response to Editor Point 2: Research Objectives}

\editortext{``There needs to be an explicit research objective(s) stated,
preferably as a separate section.''}

\textbf{Addressed.} A standalone Section~3 (``Research Objectives'') now
presents three explicit objectives:
\begin{description}[nosep]
  \item[RO1] Quantify agreement between heuristic, LLM-ensemble, and human
    judges on tool-using agent outputs.
  \item[RO2] Characterize parameter-level error propagation through agent
    pipelines under validated evaluation.
  \item[RO3] Evaluate a runtime interceptor under a concurrent control.
\end{description}
Each objective is accompanied by a rationale paragraph.

\newtext{See Section~3 of the revised manuscript.}

\section{Response to Editor Point 3: Dataset Description}

\editortext{``I have questions about your dataset. So, you need to
explicitly provide a description of your dataset(s), in the manuscript,
which is currently lacking.''}

\textbf{Addressed.} Section~4 (``Dataset: AgentProp-Bench'') now provides
a full description including:
\begin{itemize}[nosep]
  \item Task design and construction methodology (Section~4.1)
  \item Domain breakdown with Table~1 showing 2,000 tasks across
        calendar (667), weather (674), medical (559), and knowledge (100)
        domains, plus 300 held-out retail tasks (Section~4.2)
  \item Tool schema and simulated environment (Section~4.3)
  \item Error-injection protocol with a worked example trace (Section~4.4)
  \item Trace collection details: 2,300 traces across 9 models (Section~4.5)
  \item Human label collection: 100 stratified labels, annotator description,
        sampling rationale (Section~4.6)
  \item Data availability with repository structure (Section~4.7)
\end{itemize}

\newtext{See Section~4 and Table~1 of the revised manuscript.}

\section{Response to Editor Point 4: Discussion and Implications}

\editortext{``You must more clearly highlight your research's theoretical
and practical implications; it is unclear how this research differs from
existing work. There should be an explicit section addressing the
discussion of results and implications.''}

\textbf{Addressed.} Section~7 (``Discussion and Implications'') is now a
dedicated section with four subsections:
\begin{itemize}[nosep]
  \item Section~7.1: Theoretical implications---rejection and recovery as
    orthogonal axes of agent robustness; why heuristic evaluation
    systematically misranks models.
  \item Section~7.2: Practical implications for agent designers---schema
    validators, pre-execution semantic checks, interceptor applicability.
  \item Section~7.3: Practical implications for IR/IS evaluation---substring
    heuristics are chance-level, LLM-ensemble judges need validation,
    methodology transfers to RAG/QA/agent benchmarks.
  \item Section~7.4: Explicit differentiation from AgentHallu
    (Liu et al., 2026), $\tau$-bench, AgentBench, and BFCL.
\end{itemize}

\newtext{See Section~7 of the revised manuscript.}

\section{Response to Editor Point 5: Copy Editing}

\editortext{``The manuscript requires thorough copy editing. Please keep
in mind that many of the IP\&M audience members are not native English
speakers, so the manuscript must be accessible to them.''}

\textbf{Addressed.} The manuscript has been thoroughly copy-edited with
the following guidelines applied:
\begin{itemize}[nosep]
  \item Sentences longer than 25 words were split.
  \item Idioms and metaphors were replaced with direct language.
  \item Every acronym (LLM, EPS, IR, etc.) is defined on first use.
  \item Active voice is used where possible.
  \item Each paragraph addresses one idea in 3--6 sentences.
\end{itemize}

\section{Response to Editor Point 6: Manuscript Length}

\editortext{``The manuscript is short, so you have space to go into more
detail concerning this research's prior literature, methods, data, and
implications.''}

\textbf{Addressed.} The manuscript has been expanded from approximately
5,200 words to approximately 9,500 words. Major additions include:
\begin{itemize}[nosep]
  \item Expanded Related Work (Section~2): 6 thematic subsections with
        18 new references.
  \item New standalone sections: Research Objectives (Section~3), Dataset
        (Section~4), Discussion (Section~7), Limitations (Section~8).
  \item New results: Rejection vs.\ Recovery decomposition (Section~6.3),
        human-validated judge reliability (Section~6.1), concurrent-control
        interceptor evaluation (Section~6.4).
\end{itemize}

\section{Response to Editor Point 7: References and Literature Review}

\editortext{``Your references need to tie this research into the computing
and information science community \ldots\ you must update your literature
review.''}

\textbf{Addressed.} The Related Work section (Section~2) has been
reorganized into six thematic subsections and expanded with 18 new
2024--2026 citations. These include:
\begin{itemize}[nosep]
  \item IR/IS venues: at least 3 papers from SIGIR, CIKM, TOIS, and IP\&M.
  \item Agent-hallucination work: AgentHallu (Liu et al., 2026) with an
        explicit differentiation paragraph (Section~2.3).
  \item LLM-as-judge methodology: recent work on judge reliability in
        IR evaluation contexts.
  \item Runtime monitoring and self-consistency approaches.
\end{itemize}

All references follow APA formatting via \texttt{elsarticle-num-names.bst}.

\newtext{See Section~2 and the updated \texttt{refs.bib}.}

\section{Response to Editor Point 8: Author Guidelines and IP\&M Style}

\editortext{``Please review recent IP\&M articles on similar topics to
get some insights on manuscript structure, presentation, and level of
research.''}

\textbf{Addressed.} We reviewed recent IP\&M articles on LLM evaluation,
information retrieval, and agent benchmarking. The revised manuscript
adopts:
\begin{itemize}[nosep]
  \item The \texttt{elsarticle} document class per IP\&M author guidelines.
  \item A section structure consistent with empirical evaluation papers
        in IP\&M (Introduction, Related Work, Objectives, Dataset,
        Methodology, Results, Discussion, Limitations, Conclusion).
  \item APA-style references.
  \item Explicit research objectives and a dedicated limitations section.
\end{itemize}

\section{Additional Methodological Corrections Identified During Revision}

During revision we identified that the original stage-indicator
definitions in the submitted manuscript coupled $S_2$ (tool execution
with injected parameters) and $S_3$ (final answer wrong), producing a
definitionally forced $r_{2,3} = 1.00$. Specifically, the
\texttt{classify\_stage()} function in the analysis code mapped any row
with \texttt{reached\_output=True} to stage~3, and then computed $S_2$
as \texttt{stage $\geq$ 2}, which included those same rows. This made
$p_{S_2} = p_{S_3}$ for all models by construction.

We corrected the indicators to be genuinely independent: $S_2$ uses only
the \texttt{eps} column ($\text{eps} \geq 2$), and $S_3$ uses only the
\texttt{final\_correct} column. Under the corrected indicators, $p_{S_2}
\neq p_{S_3}$ for most models (e.g., DeepSeek-V3: $p_{S_2}=7.0\%$ vs.\
$p_{S_3}=29.5\%$), and $r_{2,3}$ ranges from 0.43 to 0.97 across models
(human-calibrated mean $\approx 0.62$).

Additionally, we re-judged all 1,350 P2 semantic-wrong traces with the
validated 3-LLM ensemble judge (rather than the substring heuristic used
in the original submission), collected 100 stratified human labels to
validate the ensemble at $\kappa=0.432$ vs.\ humans, and re-ran the
runtime interceptor under a concurrent $n{=}600$ no-interceptor control
(rather than comparing against a historical baseline from a different
task pool).

All numerical claims in the revised manuscript reflect the corrected
analysis. See Section~5.2 for the revised methodology and the repository's
\texttt{CORRECTED\_CLAIMS\_DIFF.md} for a one-to-one mapping between
original and corrected numbers.

We also renamed the benchmark from AgentHallu-Bench to AgentProp-Bench
to reflect its focus on propagation measurement and to avoid confusion
with the concurrently published ``AgentHallu'' paper
(Liu et al., 2026, arXiv:2601.06818), which does step-level attribution
on a different dataset.

\section{Summary of Changes}

\begin{table}[h]
\centering\small
\begin{tabular}{lcc}
\toprule
Metric & Original & Revised \\
\midrule
Word count & $\sim$5,200 & $\sim$9,500 \\
References & 28 & 46+ \\
New 2024--2026 refs & 0 & 18 \\
Sections & 5 & 9 \\
Tables & 6 & 8 + supplementary \\
Figures & 1 & 5 + supplementary \\
Human labels & 0 & 100 \\
\bottomrule
\end{tabular}
\end{table}

\end{document}
